\paragraph{Threshold stability of the matched certificates.}
We kept the adaptive branches, equal-weight batch mixtures, and adverse-record
distributions fixed and calculated a rational interval $[L,U)$ for each of
the 222 matched certificates. Throughout this interval, adaptation guarantees
success and the optimal randomized batch has worst-record success exactly
$1/2$. The lower endpoint combines the adaptive branch bound with
the two-batch covering bound; the upper endpoint is the first threshold at
which the saved adverse distribution permits success above $1/2$ for a batch.
All 222 intervals contain $1/2$ in their interior. The complete bands
$[0.49,0.51]$, $[0.475,0.525]$, and $[0.45,0.55]$ are certified on 179, 119,
and 53 blocks, respectively. The symmetric margin around $1/2$ has median
$0.0263$, interquartile range $[0.0134,0.0488]$, and minimum $0.00020$.
These are sufficient ranges for the frozen certificates; outside them this
check is inconclusive. Supplement S1 provides the exact endpoints and a
solver-free checker that reconstructs all branch and covering inequalities.

\begin{table}[htbp]
\centering
\small
\caption{Blocks retaining adaptive certainty and an exact randomized-batch
success optimum of $1/2$ at each precision threshold. The same policies and
adverse distributions are used throughout.}
\label{tab:threshold-stability}
\begin{tabular}{lrrrrrrr}
\toprule
Threshold $\tau$ & $0.45$ & $0.475$ & $0.49$ & $0.50$ & $0.51$ & $0.525$ & $0.55$\\
\midrule
Original ($n=165$) & 73 & 120 & 153 & 165 & 148 & 125 & 91\\
External ($n=57$) & 19 & 30 & 43 & 57 & 56 & 51 & 35\\
Total ($n=222$) & 92 & 150 & 196 & 222 & 204 & 176 & 126\\
\bottomrule
\end{tabular}
\end{table}
